\uuid{2oww}
\titre{Continuité en un point}
\auteur{bodin ; hanani ; mzari}
\organisation{exo7}
\licence{CC BY-NC-SA 4.0}
\datecreate{2020-01-20}
\niveau{L1}
\module{Analyse}
\chapitre{Continuité, limite et étude de fonctions réelles}
\sousChapitre{Continuité : théorie}
\difficulte{3}
\theme{Fonctions continues, continuité en un point}
\exo7qcmid{lille-1.continuite.017}

\contenu{
\begin{question}
\qclassification{123.01,123.02}
\qtags{motcle=continuité en un point}
Parmi les propriétés suivantes, quelles sont celles qui impliquent que $f$ est continue en $x_ 0$ ?
\begin{answers}
    \bad{$f(x)^2 \to f(x_0)^2$ (lorsque $x \to x_0$)}

    \good{$f(x)^3 \to f(x_0)^3$ (lorsque $x \to x_0$)}

    \bad{$E(f(x)) \to E(f(x_0))$ (lorsque $x \to x_0$)}

    \good{$\exp(f(x)) \to \exp(f(x_0))$ (lorsque $x \to x_0$)}
\end{answers}
\begin{explanations}
Si $f(x)^2 \to f(x_0)^2$  alors ce n'est pas toujours vrai que $f(x) \to f(x_0)$, prendre la fonction $f(x)=-1$ si $x<0$ et $f(x)=+1$ sinon. Par contre avec le cube c'est vrai, car la fonction $x \mapsto x^3$ est une bijection continue de $\Rr$ dans $\Rr$, idem avec l'exponentielle !
\end{explanations}
\end{question}
}
